\documentclass[11pt]{article}
\usepackage[T1]{fontenc}
\usepackage[french]{babel}
\usepackage[margin=1in]{geometry}
\usepackage{amsmath}    

\begin{document}
\title{TIPE 2026 : Accouplement magnétique et propulsion marine}
\author{Pacôme Daval \and Louis-Antoine Wen \and Chloé Blin}
\date{\today}
\maketitle

\newpage

\tableofcontents
\newpage

\section{Introduction}


\section{Modélisaton}
Le principal objectif est tout d'abord de trouver le couple maximale transmissble à travers ce type d'accouplement.
Une première approche permet de chercher les paramètres principaux qui vont influencer le couple maximal. Voici une liste les présentant:
\begin{enumerate}
  \item La distance entre les deux parties (menante et menée)
  \item La puissance des aimants, et donc leur champ magnétique rémanent
  \item Le nombre de paire de pôles
  \item La géométrie des aimants (cylindrique, annulaire, etc.)
  \item La nature du matériau séparant les deux parties
\end{enumerate}

Pour notre première expérience, nous avons choisi de faire un montage assez simple mais nous permettant de tester
un maximum de paramètres. 
\end{document}